\documentclass{article}

\usepackage{fullpage}
\usepackage{url}
\usepackage{graphicx}
\usepackage{amsmath}
\usepackage{physics}
\usepackage{mathtools}
\usepackage{bbold}

\DeclareMathOperator{\erfi}{erfi}
\newcommand{\R}{\ensuremath{\mathbb{R}}}

\title{Quantum Mechanics II Excercise 3}
\author{Idan Stark}
\date{\today}

\begin{document}
\maketitle

\section{$SU(2), SO(3)$}
\subsection{$SU(2)$ and the Pauli matrices}
$SU(2)$ is defined as the set of all $2\times2$ complex matrices such that $U^\dagger U = \mathbb{1}$ and $\det U = 1$. We can characterise those matrices by starting from a generic $2\times2$ matrix and narrowing down using the conditions.
\newline
Given a $2\times2$ matrix U, it can be written as
\begin{equation}
    U = \begin{pmatrix}a & b \\ c & d\end{pmatrix}
\end{equation}
The first condition for it to be in $SU(2)$ is $U^\dagger U = \mathbb{1}$. This means that the matrix is invertible and that its inverse is its complex conjugate
\begin{equation}
    U^\dagger = U^{-1}
\end{equation}
Using the formula for the inverse of a $2\times2$ matrix:
\begin{equation}
    U^{-1} = \frac{1}{|\det U|}\begin{pmatrix} d & -b \\ -c & a \end{pmatrix}
\end{equation}
And using the second condition $\det U = 1$ we get:
\begin{equation}
    \begin{pmatrix} a^* & c^* \\ b^* & d^* \end{pmatrix}
    =
    \begin{pmatrix} d & -b \\ -c & a \end{pmatrix}
\end{equation}
\begin{equation}
    a^* = d \qquad d^* = a \qquad b^* = -c \qquad c^* = -b
\end{equation}
In other words, we can write $U$ as:
\begin{equation}
    U = \begin{pmatrix} a & b \\ -b^* & a^* \end{pmatrix}
\end{equation}
Writing $a$ and $b$ in their complex cartesian representation, we get:
\begin{equation}
    \label{eq:final_u}
    U = \begin{pmatrix} x + iy & w + iz \\ -w + iz & x - iy \end{pmatrix}
\end{equation}
And requiring $\det U = 1$ gives us:
\begin{equation}
    \label{eq:final_u_condition}
    x^2 + y^2 + w^2 + z^2 = 1
\end{equation}
Now, let us look at the expression
\begin{equation}
    A = \exp\left[\frac{1}{2} i\theta_i \sigma_i\right]
\end{equation}
Using the definition of the Pauli matrices, we get
\begin{equation}
    A = \exp\left[\frac{1}{2}i\theta_i \sigma_i\right] = \exp
        \begin{pmatrix}
            i\theta_z & \theta_y + i\theta_x \\
            - \theta_y+ i\theta_x & -i\theta_z
        \end{pmatrix} \equiv \exp P(\theta)
\end{equation}
Where we named the matrix in the exponent $P(\theta)$. We can now expand the exponent, but first let us calculate the powers of $P(\theta)$:
\begin{equation}
    P^2(\theta) = \begin{pmatrix}
            i\theta_z & \theta_y + i\theta_x \\
            - \theta_y+ i\theta_x & -i\theta_z
        \end{pmatrix}\begin{pmatrix}
            i\theta_z & \theta_y + i\theta_x \\
            - \theta_y+ i\theta_x & -i\theta_z
        \end{pmatrix}
        =
        -\begin{pmatrix}
            \theta_x^2 + \theta_y^2 + \theta_z^2 & 0 \\ 0 & \theta_x^2 + \theta_y^2 + \theta_z^2
        \end{pmatrix} \equiv -\theta^2 \mathbb{1}
\end{equation}
Therefore, we can split the powers of $P$ to even and odd powers:
\begin{equation}
    P^{2n}(\theta) = \left(-1\right)^n\theta^{2n}\mathbb{1} \qquad P^{2n+1}(\theta) = \left(-1\right)^n\theta^{2n}P
\end{equation}
And thus
\begin{equation}
\begin{gathered}
    A = \exp P = \sum_{n=0}^\infty \frac{P^n}{n!} = \mathbb{1}\sum_{n \ even}\frac{\left(i\theta\right)^n}{n!} + P\sum_{n \ odd}\frac{\left(i\theta\right)^{n - 1}}{n!} = \cos(\theta)\mathbb{1} + \sin(\theta) \frac{P}{\theta} = \\
    = \begin{pmatrix}
        \cos(\theta) + i\sin(\theta)\frac{\theta_z}{\theta} & \sin(\theta)\frac{\theta_y + i\theta_x}{\theta} \\ 
        \sin(\theta)\frac{-\theta_y + i\theta_x}{\theta} & \cos(\theta) - i\sin(\theta)\frac{\theta_z}{\theta}
    \end{pmatrix}
\end{gathered}
\end{equation}
The determinant of this matrix is:
\begin{equation}
\begin{gathered}
    \det A = \left(\cos(\theta) + i\sin(\theta)\frac{\theta_z}{\theta}\right) \left(\cos(\theta) - i\sin(\theta)\frac{\theta_z}{\theta}\right) - \sin^2(\theta)\frac{- \theta_x^2 - \theta_y^2}{\theta^2} = 
    \\
    = \cos^2(\theta) + \sin^2(\theta)\frac{\theta_z^2}{\theta^2} + \sin^2(\theta)\frac{\theta_x^2 + \theta_y^2}{\theta^2} = \cos^2(\theta) + \sin^2(\theta)\frac{\theta_x^2 + \theta_y^2 + \theta_z^2}{\theta^2} =
    \\
    = \cos^2(\theta) + \sin^2(\theta) = 1
\end{gathered} 
\end{equation}
Which means that with a small change of variables we can make it fit the matrix in equation \ref{eq:final_u}:
\begin{equation}
    x = \cos(\theta) \qquad y = \sin(\theta)\frac{\theta_z}{\theta}
    \qquad w = \sin(\theta)\frac{\theta_y}{\theta} \qquad z = \sin(\theta)\frac{\theta_x}{\theta}
\end{equation}
Which means that for every $2\times2$ matrix $U$ in $SU(2)$ one can find $\theta_x, \theta_y, \theta_z$ such that $U = \exp\left[\frac{1}{2}i \theta_i\sigma_i\right]$, and thus $t_i = \frac{1}{2}\sigma_i$ furnish a representation of the $su(2)$ algebra. 
\subsection{The $SO(3)$ group}
For the $SO(3)$ group, we can follow the same pattern. Starting from a $3\times3$ matrix
\begin{equation}
    O = \begin{pmatrix}
        a & b & c \\
        d & e & f \\
        h & k & l
    \end{pmatrix}
\end{equation}
And demanding that $O^TO = \mathbb{1}$, we get:
\begin{equation}
    O^T = \begin{pmatrix}
        a & d & h \\
        b & e & k \\
        c & f & l
    \end{pmatrix}
\end{equation}
\begin{equation}
    O^TO = \begin{pmatrix}
        a^2 + d^2 + h^2 & ab + de + hk & ac + df + hl \\
        ab + de + hk & b^2 + e^2 + k^2 & bc + ef + kl \\
        ac + df + hl & bc + ef + kl & c^2 + f^2 + l^2
    \end{pmatrix}
\end{equation}
Which gives the equations
\begin{equation}
    a^2 + d^2 + h^2 = b^2 + e^2 + k^2 = c^2 + f^2 + l^2 = 1
\end{equation}
\begin{equation}
    ab + de + hk = ac + df + hl = bc + ef + kl = 0
\end{equation}
The three matrices that make up the representation of the algebra are:
\begin{equation}
    L_1 = \begin{pmatrix}
        0 & 0 & 0 \\
        0 & 0 & -i \\
        0 & i & 0
    \end{pmatrix}
    \qquad
    L_2 = \begin{pmatrix}
        0 & 0 & i \\
        0 & 0 & 0 \\
        -i & 0 & 0
    \end{pmatrix}
    \qquad
    L_3 = \begin{pmatrix}
        0 & -i & 0 \\
        i & 0 & 0 \\
        0 & 0 & 0
    \end{pmatrix}
\end{equation}
Their linear sum with constant $\varphi_i$ is
\begin{equation}
    P = (-i) i \begin{pmatrix}
        0 & -\varphi_3 & \varphi_2 \\
        \varphi_3 & 0 & -\varphi_1 \\
        -\varphi_2 & \varphi_1 & 0
    \end{pmatrix} = \begin{pmatrix}
        0 & -\varphi_3 & \varphi_2 \\
        \varphi_3 & 0 & -\varphi_1 \\
        -\varphi_2 & \varphi_1 & 0
    \end{pmatrix}
\end{equation}
Which means that
\begin{equation}
    A = \exp\left[-i\varphi_i L_i\right] = \exp P = \exp \begin{pmatrix}
        0 & -\varphi_3 & \varphi_2 \\
        \varphi_3 & 0 & -\varphi_1 \\
        -\varphi_2 & \varphi_1 & 0
    \end{pmatrix}
\end{equation}
We can now look at the powers of $P$ to get
\begin{equation}
    P^2 = \begin{pmatrix}
        0 & -\varphi_3 & \varphi_2 \\
        \varphi_3 & 0 & -\varphi_1 \\
        -\varphi_2 & \varphi_1 & 0
    \end{pmatrix}\begin{pmatrix}
        0 & -\varphi_3 & \varphi_2 \\
        \varphi_3 & 0 & -\varphi_1 \\
        -\varphi_2 & \varphi_1 & 0
    \end{pmatrix}
    = \begin{pmatrix}
        -\varphi_2^2 - \varphi_3^2 & \varphi_1\varphi_2 & \varphi_1\varphi_3 \\ \varphi_1\varphi_2 & -\varphi_1^2 - \varphi_3^2 & \varphi_2\varphi_3 \\ \varphi_1\varphi_3 & \varphi_2\varphi_3 & -\varphi_1^2 - \varphi_2^2
    \end{pmatrix}
\end{equation}
\begin{equation}
    P^3 = \begin{pmatrix}
        0 & -\varphi_3 & \varphi_2 \\
        \varphi_3 & 0 & -\varphi_1 \\
        -\varphi_2 & \varphi_1 & 0
    \end{pmatrix} \begin{pmatrix}
        -\varphi_2^2 - \varphi_3^2 & \varphi_1\varphi_2 & \varphi_1\varphi_3 \\ \varphi_1\varphi_2 & -\varphi_1^2 - \varphi_3^2 & \varphi_2\varphi_3 \\ \varphi_1\varphi_3 & \varphi_2\varphi_3 & -\varphi_1^2 - \varphi_2^2
    \end{pmatrix} = -\varphi^2 P
\end{equation}
Then, for $n > 0$
\begin{equation}
    P^{2n} = (i\varphi)^{2n-2} P^{2} \qquad P^{2n+1} = (i\varphi)^{2n}P
\end{equation}
\begin{equation}
\begin{gathered}
    A = \exp P = \sum_{n=0}^{\infty} \frac{P^n}{n!} =
    \sum_{n \ even}\frac{P^n}{n!} + \sum_{n \ odd}\frac{P^n}{n!} = 
    \mathbb{1} + \sum_{n > 0 \ even}\frac{(i\varphi)^{n - 1}}{n!}P^2 + \sum_{n \ odd}\frac{(i\varphi)^n}{n!}P = 
    \\
    = \mathbb{1} - i \frac{\cos(\varphi) - 1}{\varphi}P^2 + \sin(\varphi)P
\end{gathered}
\end{equation}
\subsection{Commutation relations}
It's easy to see that for $t_i = \frac{1}{2}\sigma_i$, $\left[t_i,t_j\right] = i\varepsilon_{ijk}t_k$. The simplest way is just to check them one by one:
\begin{equation}
\begin{gathered}
    \left[t_1,t_1\right] = 0
    \qquad
    \left[t_2,t_2\right] = 0
    \qquad
    \left[t_3,t_3\right] = 0
    \\
    \left[t_1,t_2\right] = - \left[t_2,t_1\right] = t_1t_2 - t_2t_1 = \\ = \frac{1}{4}\left[\begin{pmatrix}0&1\\1&0\end{pmatrix}\begin{pmatrix}0&-i\\i&0\end{pmatrix} - \begin{pmatrix}0&-i\\i&0\end{pmatrix}\begin{pmatrix}0&1\\1&0\end{pmatrix}\right] = \frac{1}{4}\left[\begin{pmatrix}i&0\\0&-i\end{pmatrix} - \begin{pmatrix}-i&0\\0&i\end{pmatrix}\right] = \\ = \frac{1}{4}\left[i\sigma_3 - \left(-i\right)\sigma_3\right] = \frac{1}{2}i\sigma_3 = it_3
    \\
    \left[t_1,t_3\right] = - \left[t_3,t_1\right] = t_1t_3 - t_3t_1 = \\ = \frac{1}{4}\left[\begin{pmatrix}0&1\\1&0\end{pmatrix}\begin{pmatrix}1&0\\0&-1\end{pmatrix} - \begin{pmatrix}1&0\\0&-1\end{pmatrix}\begin{pmatrix}0&1\\1&0\end{pmatrix}\right] = \frac{1}{4}\left[\begin{pmatrix}0&-1\\1&0\end{pmatrix} - \begin{pmatrix}0&1\\-1&0\end{pmatrix}\right] = \\ = \frac{1}{4}\left[-i\sigma_2 + \left(-i\right)\sigma_2\right] = -\frac{1}{2}i\sigma_2 = -it_2
    \\
    \left[t_2,t_3\right] = - \left[t_3,t_2\right] = t_2t_3 - t_3t_2 = \\ = \frac{1}{4}\left[\begin{pmatrix}0&-i\\i&0\end{pmatrix}\begin{pmatrix}1&0\\0&-1\end{pmatrix} - \begin{pmatrix}1&0\\0&-1\end{pmatrix}\begin{pmatrix}0&-i\\i&0\end{pmatrix}\right] = \frac{1}{4}\left[\begin{pmatrix}0&i\\i&0\end{pmatrix} - \begin{pmatrix}0&-i\\-i&0\end{pmatrix}\right] = \\ = \frac{1}{4}\left[i\sigma_1 - \left(-i\right)\sigma_1\right] = \frac{1}{2}i\sigma_1 = it_1
\end{gathered}
\end{equation}
Similarly, for $Li$ we get
\begin{equation}
\begin{gathered}
    \left[L_1,L_1\right] = 0
    \qquad
    \left[L_2,L_2\right] = 0
    \qquad
    \left[L_3,L_3\right] = 0
    \\
    \left[L_1,L_2\right] = - \left[L_2,L_1\right] = L_1L_2 - L_2L_1 = \\ = 
    \begin{pmatrix} 0 & 0 & 0 \\ 0 & 0 & -i \\ 0 & i & 0 \end{pmatrix}
    \begin{pmatrix} 0 & 0 & i \\ 0 & 0 & 0 \\ -i & 0 & 0 \end{pmatrix}
    -
    \begin{pmatrix} 0 & 0 & i \\ 0 & 0 & 0 \\ -i & 0 & 0 \end{pmatrix}
    \begin{pmatrix} 0 & 0 & 0 \\ 0 & 0 & -i \\ 0 & i & 0 \end{pmatrix}
    = \\ =
    \begin{pmatrix} 0 & 0 & 0 \\ -1 & 0 & 0 \\ 0 & 0 & 0 \end{pmatrix}
    -
    \begin{pmatrix} 0 & -1 & 0 \\ 0 & 0 & 0 \\ 0 & 0 & 0 \\ \end{pmatrix}
    = i\begin{pmatrix} 0 & -i & 0 \\ i & 0 & 0 \\ 0 & 0 & 0\end{pmatrix} = iL_3
    \\
    \left[L_1,L_3\right] = - \left[L_3,L_1\right] = L_1L_3 - L_3L_1 = \\ = 
    \begin{pmatrix} 0 & 0 & 0 \\ 0 & 0 & -i \\ 0 & i & 0 \end{pmatrix}
    \begin{pmatrix} 0 & -i & 0 \\ i & 0 & 0 \\ 0 & 0 & 0\end{pmatrix}
    -
    \begin{pmatrix} 0 & -i & 0 \\ i & 0 & 0 \\ 0 & 0 & 0\end{pmatrix}
    \begin{pmatrix} 0 & 0 & 0 \\ 0 & 0 & -i \\ 0 & i & 0 \end{pmatrix}
    = \\ =
    \begin{pmatrix} 0 & 0 & 0 \\ 0 & 0 & 0 \\ -1 & 0 & 0 \end{pmatrix}
    -
    \begin{pmatrix} 0 & 0 & -1 \\ 0 & 0 & 0 \\ 0 & 0 & 0 \\ \end{pmatrix}
    = i\begin{pmatrix} 0 & 0 & -i \\ 0 & 0 & 0 \\ i & 0 & 0\end{pmatrix} = -iL_2
    \\
    \left[L_2,L_3\right] = - \left[L_3,L_2\right] = L_2L_3 - L_3L_2 = \\ = 
    \begin{pmatrix} 0 & 0 & i \\ 0 & 0 & 0 \\ -i & 0 & 0\end{pmatrix}
    \begin{pmatrix} 0 & -i & 0 \\ i & 0 & 0 \\ 0 & 0 & 0\end{pmatrix}
    -
    \begin{pmatrix} 0 & -i & 0 \\ i & 0 & 0 \\ 0 & 0 & 0\end{pmatrix}
    \begin{pmatrix} 0 & 0 & i \\ 0 & 0 & 0 \\ -i & 0 & 0\end{pmatrix}
    = \\ =
    \begin{pmatrix} 0 & 0 & 0 \\ 0 & 0 & 0 \\ 0 & -1 & 0 \end{pmatrix}
    -
    \begin{pmatrix} 0 & 0 & 0 \\ 0 & 0 & -1 \\ 0 & 0 & 0 \\ \end{pmatrix}
    = i\begin{pmatrix} 0 & 0 & 0 \\ 0 & 0 & -i \\ 0 & i & 0\end{pmatrix} = iL_1
\end{gathered}
\end{equation}
This means that both $t_i$ and $L_i$ satisfy the same commutation relations, and there is an obvious isomorphism between their respective algebras.

\section{A bit of $SO(1,3)$}
\section{What adds under Lorentz transformations?}
\end{document}